\section{Monodromy, Stokes data, and the character variety}\label{sec:monodromy}

Sections \ref{sec:point} and \ref{sec:projective} identified System $\mathsf C$
with a member of the van der Put--Saito family at a located point of the chart
$U_0$, and removed the projective qualifier from that identification. What has
not yet been produced is the \emph{monodromy} data: the point of the character
variety that System $\mathsf C$ represents. This section produces it. It does
so with a solver that is first validated against a case with closed-form
answers, and under a set of orientation and ordering conventions that are
\emph{selected by test} rather than assumed.

The distinction matters more here than anywhere else in the paper. Trace
coordinates are not convention-independent: a different choice of loop
orientation, of the order in which local monodromies are composed, or of which
loop at infinity is called $\gamma_\infty$, produces a different triple for the
same differential system. A stated triple without a stated convention is not a
result. The convention used below is therefore the output of
\S\ref{ssec:selection}, not an input to it.

\subsection{The solver, and a case where the answer is
known}\label{ssec:kummer}

The solver computes formal solutions at the irregular singular point, the
Stokes matrices relating consecutive sectors, and the monodromy
representation, by numerical analytic continuation of a fundamental matrix. It
was validated on Kummer's equation, whose monodromy, formal data and Stokes
data are all available in closed form.

\begin{remark}[Known-answer validation; \texttt{S22-01}, \textbf{VERIFIED}]
The solver reproduces Kummer's monodromy, formal and Stokes data against the
closed forms: $32$ of $32$ checks pass at at least $35$ digits, with $44$ to
$151$ digits observed, at both pre-registered parameter values.
\end{remark}

Two structural facts about Kummer's case were established in the course of that
validation, and both are used later.

\begin{remark}[The trace-free gauge is not optional; \texttt{S22-02},
\textbf{STRUCTURAL}]
Kummer's companion system is not trace-free: $\operatorname{tr}B=1-b/z$. Hence
$\det M_0=e^{-2\pi i b}\neq 1$, and the van der Put--Saito $\mathrm{SL}_2$
factorisation is unavailable until the trace-free gauge of
\S\ref{ssec:tracefree} has been applied. This is the same reduction that
\S\ref{sec:projective} performed for System $\mathsf C$, met here in a case
where its necessity is visible from the local exponents alone.
\end{remark}

\begin{remark}[Orientation at infinity; \texttt{S22-03}, \textbf{STRUCTURAL}]
The counterclockwise circle at infinity is homotopic to $\gamma_0$. Van der
Put--Saito's $M_\infty$ is the \emph{clockwise} loop, and their $\alpha$ is
consequently $e^{-2\pi i L_{11}}$. The sign in the exponent is the whole
content of the remark: it is the difference between two conventions that agree
on every other datum, and it is the reason \S\ref{ssec:selection} tests
orientation rather than fixing it.
\end{remark}

Two further algebraic checks from the same validation are omitted from this
section deliberately. Both were \emph{vacuous as first written} --- one an
identity that holds for every unimodular matrix with a nonzero corner entry,
the other a comparison across two different factorisation orderings --- and
both were repaired and proved. They record how the validation battery was
debugged, not what the battery established, and are set out in the methods
appendix, \S\ref{app:methods}, with the claims \texttt{S22-04} and
\texttt{S22-05}. The corrected form of the second is in any case a
\emph{dependent} restatement of checks already counted, and is not independent
evidence.

\subsection{The system at the located point}\label{ssec:atpoint}

Before any monodromy is computed, the family member at the located point is
checked against the local data it is supposed to have.

\begin{remark}[The family at the point; \texttt{S22-06}, \textbf{VERIFIED}]
At the point located in \S\ref{sec:point}, the printed $(4.1)$ family satisfies
$\det A_0=\det A_1=0$ exactly, and $A_\infty=(t/2)\operatorname{diag}(-1,1)$ to
$121$ digits.
\end{remark}

The first of these is exact rather than numerical because it is a polynomial
identity in the entries; the second is the residue at infinity, and is the
datum the whole Riemann--Hilbert side depends on. One discrepancy with the
printed source surfaced at this stage and is recorded here because a reader
checking against that source will meet it.

\begin{remark}[A discrepancy with the printed prose; \texttt{S22-07},
\textbf{VERIFIED}]
The computation gives $(A_0+A_1)_{11}=-\theta_\infty/2$, where the prose of van
der Put--Saito at p.~2629 states $\theta_\infty$. This was first found in the
location analysis and is reconfirmed here at the numeric point. \emph{Nothing
in this paper rests on it}: the factor is never used, and it is reported so
that the discrepancy is on record rather than rediscovered.
\end{remark}

\begin{remark}[The exponent at infinity, in closed form; \texttt{S22-08},
\textbf{VERIFIED}]
$L_{11}=\theta_\infty/2$ exactly, so that $\alpha=e^{\pm i\pi\theta_\infty}$ in
closed form, with $|\alpha|=1$ and
$|\arg\alpha|=\pi\sqrt3/9=0.60459978807807261686$. Both the modulus and the
argument were predicted before the run, and both were hit.
\end{remark}

That $|\alpha|=1$ is worth a sentence. It says the monodromy eigenvalue at
infinity lies on the unit circle, which is what one expects of a $\theta_\infty$
that is real --- and $\theta_\infty=-2\pm\sqrt3/9$ is real, by
\eqref{eq:eliminant}. The prediction was registered in advance precisely
because a modulus differing from $1$ would have indicated a transcription error
in the exponent rather than a feature of the system.

The formal series at the irregular point is divergent, so the accuracy
available from it is bounded by optimal truncation, and that bound has to be
measured rather than assumed.

\begin{remark}[Optimal truncation; \texttt{S22-09}, \textbf{VERIFIED}]
The truncated formal series' residual sits at the minimal term of the series
--- ratio $1.29$ against a bound of $N^3$ --- and decays at the predicted
Gevrey rate $e^{-|t|R}$ between $R=160$ and $R=140$, with observed over
predicted $1.259$.
\end{remark}

This is the accuracy ceiling for everything downstream. It is not a limitation
of the arithmetic precision, and \S\ref{ssec:notroundoff} shows directly that
raising the precision does not move it.

\subsection{Selecting the convention}\label{ssec:selection}

The conventions were not chosen. A candidate set was enumerated, four
conditions were registered in advance, and the set was searched.

The four conditions are: that the monodromy representation act on the right;
that the small loops around the finite punctures run counterclockwise; that
$M_\infty$ be the clockwise loop at infinity, as \texttt{S22-03} requires; and
that van der Put--Saito's $M_1$ be the loop $\gamma_1$. The test is the
relation $M_0M_1M_\infty=I$, at a threshold of $10^{-30}$.

\begin{remark}[The candidate set was too small; \texttt{S22-10},
\textbf{STRUCTURAL}]
The pre-registered $8$-element candidate set tied the orientation of
$\gamma_\infty$ to that of the small loops, and therefore could not represent
$\gamma_0\gamma_1=\gamma_\infty$ at all. None of the $8$ passed; the best
residual was $\|MMM-I\|=6.85$. The set was widened to $16$ by untying the two
orientations, with the four conditions and the $10^{-30}$ threshold left
unchanged.
\end{remark}

The widening is recorded rather than quietly performed because it is a change
to a pre-registered object. What was \emph{not} changed is the criterion: the
same four conditions and the same threshold were applied to the larger set. A
failure of all eight candidates at $O(1)$ is not a marginal miss, and it is the
signature of a set that omits the truth rather than of a threshold set too
tight.

\begin{remark}[Exactly one convention survives; \texttt{S22-11},
\textbf{VERIFIED}]
Over the $16$ candidates exactly one satisfies all four conditions: right
action, small loops counterclockwise, $M_\infty$ clockwise, van der
Put--Saito's $M_1=\gamma_1$. For it, $\|M_0M_1M_\infty-I\|=2.44\times10^{-105}$.
The other $15$ fail, several on $\alpha$ alone, with residuals $0.327$, $2.27$
and $2.34$.
\end{remark}

The gap between $2.44\times10^{-105}$ and the nearest competitor is not a
matter of ranking: the surviving candidate satisfies the relation to the full
working precision, and the others miss it by $O(1)$. \emph{Exactly one} is
therefore a genuine selection and not a choice among near-misses, which is what
makes the triple below well defined.

\subsection{The trace coordinates}\label{ssec:traces-value}

Throughout, $B_1,B_2,B_3$ are the monodromy matrices in the sense fixed in
\S\ref{ssec:traces}, with $a_i=\operatorname{tr}(B_i)\equiv 2$, and the affine
trace coordinates are
\begin{equation}\label{eq:traces-restate}
x_1=\operatorname{tr}(B_2B_3),\qquad
x_2=\operatorname{tr}(B_1B_3),\qquad
x_3=\operatorname{tr}(B_1B_2),
\end{equation}
as in \eqref{eq:n-traces}. We restate \eqref{eq:traces-restate} here rather
than refer to it, because the values below are meaningless without knowing
which product each coordinate is the trace of, and because $x_1$ is the trace
of the product that omits $B_1$ --- the indexing most easily got backwards.

\begin{remark}[The values; \texttt{S22-12}, \textbf{VERIFIED}]
At $t_0$, in the convention selected in \S\ref{ssec:selection},
\begin{align*}
x_1&=-3.4791908904110720783362611334763347
 -1.085013570276491408740290278433511\,i,\\
x_2&=\hphantom{-}5.0274753770206575050513601616191382
 +2.3656209874946766465307178962989066\,i,\\
x_3&=\hphantom{-}1.9725752524636335304002661101339216,\\
\alpha&=\hphantom{-}0.82272965712005465178275844880452679
 +0.56843285557321306726020179279562511\,i.
\end{align*}
The error bar is $8.49\times10^{-105}$, obtained from step halving --- $273$
steps against $487$ --- and not from a single
run.
\end{remark}

That the error estimate comes from step halving rather than from a residual
reported by a single integration is the point of the parenthesis. A single run
can report a small residual and be wrong by far more; two runs at different
step sizes bound the discretisation error they do not share. Note also that
$x_3$ is real while $x_1$ and $x_2$ are not, and that
$|\alpha|=1$ to the digits shown, consistent with \texttt{S22-08}.

\begin{remark}[A smooth point of a singular fibre; \texttt{S22-13},
\textbf{VERIFIED}]
The point is a \emph{smooth} point of a \emph{singular} cubic fibre. The fibre
is singular because $s_1=s_2=2$ places it on two resonance divisors; but
$\operatorname{grad}F$ has modulus $(1.523094,\,1.5231748,\,21.240405)$
componentwise there, which is nonvanishing, so the coordinates retain
first-order sensitivity at the point.
\end{remark}

The two adjectives pull in opposite directions and both are needed. Singularity
of the fibre is what \S\ref{ssec:resonance} warned of: the resonance conditions
are not excluded, and they are met here. Smoothness of the point is what
rescues the computation, because a coordinate system that degenerated at the
point would make the triple above uninformative however many digits it carried.

\subsection{The coordinates see the leaf}\label{ssec:sensitivity}

One thing remains to be checked before \S\ref{sec:leaf} can use these
coordinates as a test of leaf membership: that they move when the point moves
off the leaf. An invariant that is constant for a trivial reason confirms
nothing.

\begin{remark}[Sensitivity; \texttt{S22-14}, \textbf{VERIFIED}]
The observable is sensitive to leaf membership: at $\delta=10^{-6}$,
$\max_i|dx_i|=8.686\times10^{-6}$ in the $q$ direction and
$2.639\times10^{-5}$ in the $p$ direction, against a floor of $10^{-15}$, with
first-order ratios $99.99978$ and $99.99956$, both inside the pre-registered
window $[50,200]$.
\end{remark}

The ratios are the substance, not the magnitudes. They say the response is
first order in the displacement, which is what distinguishes a coordinate that
genuinely resolves the leaf direction from one that happens to be small. The
constancy reported in \S\ref{sec:leaf} is therefore a statement with content:
the same coordinates, displaced off the leaf by $10^{-6}$, move by five orders
of magnitude more than the floor.
